\documentclass{article}
\usepackage{graphicx} % Required for inserting images

\title{Álgebra - TPC4}
\author{Afonso V. Francisco - 106618 }
\date{15 Março 2024}
\usepackage{amsfonts}
\usepackage{amssymb}
\begin{document}

\maketitle

\textbf{Seja $G=C^x$ (números complexos não nulos, com a operação de multiplicação) e seja  $H = \{-1,+1,+i,-i\}$ o subgrupo de $G$ formado pelas raízes quartas de 1. Descreva explicitamente as classes esquerdas de $H$ em $G$. Será $G/H$ isomorfo a $G$?}\\

$G=\mathbb{C}^{\times}$, $H=\{-1,1,-i,i\}$.\\
As classes esquerdas de $H$ em $G$ são $zH=\{-z,z,-zi,zi\}$, para $z \in \mathbb{C}^{\times}$. $zH=\{\rho e^{i({\pi + \theta})},\rho e^{i \theta},\rho e^{i({3/2\pi + \theta})},\rho e^{i({\pi/2 + \theta})}\}=\{-a-bi,a+bi,b-ai,-b+ai\}$, onde usámos $z=\rho e^{i \theta }$ e $z=a+bi$, respetivamente.\\


Para aplicarmos o \underline{1º Teorema do isomorfismo} temos que encontrar um um homomorfismo de grupos sobrejetivo.
$$Seja \;\;\;\;\; f:G \rightarrow G,\;\;\; f(z)=z^4$$ (escolhido desta forma para que tenhamos $H=ker f$ ).

\begin{itemize}
    \item \underline{f é homomorfismo de grupos}:
    $f(zw)=(zw)^4=z^4w^4=f(z)f(w)$
    \item \underline{Núcleo de f}:
    $ker f=\{x \in G \mid f(x)=1_G\}=\{x \in G \mid f(x)=1\}=\{x \in G \mid x^4=1\}=\{-1,1,-i,i\}=H$
    \item \underline{Imagem de f}:
    $im f = G$. Basta ver que todos os elementos de $G$ são atingidos por $f$, já que para qualquer $x \in G$ temos $x^{1/4} \in G$.
    Portanto, f é sobrejetiva.
\end{itemize}

Assim como $f$ é um homomorfismo de grupos sobrejetivo, pelo 1º Teorema do isomorfismo temos que $f$ induz um isomorfismo $\overline{f}:\frac{G}{ker f} \rightarrow G$. Ou seja $\frac{G}{ker f}$ é isomorfo a $G$, mas como vimos anteriormente que $ker f=H$, então $\frac{G}{H}$ é isomorfo a $G$.
\end{document}
